% ============================================================================
% CHAPTER 9 — COMPARATIVE ANALYSIS
% ============================================================================
\chapter{Comparative Analysis}
\label{ch:comparative}

\begin{motivation}
    The previous six chapters each presented a single ML-enhanced approach
    in isolation.  This chapter pulls back to a bird's-eye view and
    compares all approaches side-by-side across multiple performance
    dimensions: PAPR reduction, BER impact, computational complexity,
    side information requirements, and practical implementability.
    The goal is to help you understand the \emph{real trade-offs}
    when choosing a technique for a specific system.
\end{motivation}


% ----------------------------------------------------------------------------
\section{Summary Comparison Table}
\label{sec:comparison_table}
% ----------------------------------------------------------------------------

Table~\ref{tab:master_comparison} provides a comprehensive side-by-side
comparison of all six ML-enhanced approaches studied in this document.

\begin{table}[H]
    \centering
    \caption{Master comparison of ML-enhanced PAPR reduction approaches.
             Data compiled from~\cite{kim2017novel,zou2021novel,ali2016new,hao2019extended,nguyen2022deep,alnaseri2025deep}.}
    \label{tab:master_comparison}
    \scriptsize
    \begin{tabularx}{\textwidth}{@{} l l *{6}{>{\centering\arraybackslash}X} @{}}
        \toprule
        & & \textbf{PRNet} & \textbf{NN\nobreakdash-SCF} & \textbf{GA\nobreakdash-SLM}
        & \textbf{ESLM\nobreakdash-AE} & \textbf{PR\nobreakdash-DUN} & \textbf{DL\nobreakdash-AE} \\
        & & \textbf{(Ch.~3)} & \textbf{(Ch.~4)} & \textbf{(Ch.~5)}
        & \textbf{(Ch.~6)} & \textbf{(Ch.~7)} & \textbf{(Ch.~8)} \\
        \midrule
        \multicolumn{2}{l}{\textbf{Citation}}
        & \cite{kim2017novel} & \cite{zou2021novel} & \cite{ali2016new}
        & \cite{hao2019extended} & \cite{nguyen2022deep} & \cite{alnaseri2025deep} \\
        \multicolumn{2}{l}{\textbf{Year}}
        & 2017 & 2021 & 2016 & 2019 & 2022 & 2025 \\
        \midrule
        \multicolumn{2}{l}{\textbf{ML Type}}
        & DNN/AE & DNN & GA & DNN/AE & Unfold & DNN/AE \\
        \multicolumn{2}{l}{\textbf{Framework}}
        & Autoencoder & SCF & SLM+GA & SLM+AE & PGD & Autoencoder \\
        \multicolumn{2}{l}{\textbf{Domain}}
        & RF OFDM & RF OFDM & RF OFDM & DCO & RF OFDM & CO-OFDM \\
        \midrule
        \multicolumn{2}{l}{\textbf{Subcarriers ($N$)}}
        & 64 & 128 & 128 & 128 & 64--256 & 64 \\
        \multicolumn{2}{l}{\textbf{Modulation}}
        & QPSK & QPSK & QPSK & 4-QAM & 16-QAM & 16-QAM \\
        \multicolumn{2}{l}{\textbf{Channel}}
        & AWGN & AWGN & N/A & Fading & AWGN & Fiber \\
        \midrule
        \multicolumn{2}{l}{\textbf{PAPR reduction}}
        & 3--4~dB & 3--5~dB & 2--4~dB & 3--4~dB & 3--5~dB & 3--4~dB \\
        \multicolumn{2}{l}{\textbf{SI required?}}
        & No & No & Yes & No & No & No \\
        \multicolumn{2}{l}{\textbf{Distortion-free?}}
        & No & No & Yes & Yes & No & No \\
        \multicolumn{2}{l}{\textbf{BER impact}}
        & Minimal & Tunable & None & Minimal & Minimal & Minimal \\
        \midrule
        \multicolumn{2}{l}{\textbf{Training}}
        & Offline & Offline & None & Offline & Offline & Offline \\
        \multicolumn{2}{l}{\textbf{Trainable params}}
        & $\sim 10^6$ & $\sim 10^6$ & 0 & $\sim 10^5$ & $\sim 20$ & $\sim 10^5$ \\
        \multicolumn{2}{l}{\textbf{Inference cost}}
        & 1 fwd pass & 2 fwd pass & $PG$ IFFTs & 1 fwd pass & $L$ layers & 1 fwd pass \\
        \bottomrule
    \end{tabularx}
\end{table}


% ----------------------------------------------------------------------------
\section{PAPR Reduction Performance}
\label{sec:comparison_papr}
% ----------------------------------------------------------------------------

All six approaches achieve roughly 3--5~dB PAPR reduction at
$\mathrm{CCDF} = 10^{-3}$, which is comparable to conventional SLM with
$U = 4$--$16$ candidates~\cite{da2024survey}.

\begin{keyinsight}
    The PAPR reduction numbers are surprisingly similar across all approaches.
    This suggests that the \emph{fundamental limit} of achievable PAPR
    reduction (without increasing average power or using reserved tones)
    may be around 5~dB for these system sizes.
    The real differentiators are complexity, SI requirements, and BER
    impact---not raw PAPR performance~\cite{da2024survey}.
\end{keyinsight}


% ----------------------------------------------------------------------------
\section{Computational Complexity}
\label{sec:comparison_complexity}
% ----------------------------------------------------------------------------

Table~\ref{tab:complexity_comparison} breaks down the per-symbol
computational cost of each approach~\cite{kim2017novel,zou2021novel,ali2016new,hao2019extended,nguyen2022deep,alnaseri2025deep}.

\begin{table}[H]
    \centering
    \caption{Computational complexity comparison (per OFDM symbol).
             ``Fwd pass'' means one forward pass through the DNN.
             Data from~\cite{kim2017novel,zou2021novel,ali2016new,nguyen2022deep}.}
    \label{tab:complexity_comparison}
    \begin{tabular}{@{} l l l @{}}
        \toprule
        \textbf{Approach} & \textbf{Inference Cost} & \textbf{Notes} \\
        \midrule
        Classical SLM & $U \cdot \mathcal{O}(N\log N)$ & $U$ IFFTs \\
        PRNet & 1 DNN forward pass & $\sim$$N^2$ multiplications \\
        NN-SCF & 2 DNN forward passes & Tx + Rx modules \\
        GA-SLM & $PG \cdot \mathcal{O}(N\log N)$ & $P \times G$ IFFTs \\
        ESLM-AE & 1 DNN fwd + 1 IFFT & Phase gen + IFFT \\
        PR-DUN & $L \times \mathcal{O}(N)$ & $L$ gradient + proj. \\
        DL-AE & 1 DNN forward pass & Similar to PRNet \\
        \bottomrule
    \end{tabular}
\end{table}

\begin{keyinsight}
    PR-DUN stands out with by far the fewest trainable parameters ($\sim 20$)
    and a per-layer cost that is only $\mathcal{O}(N)$ (a gradient computation
    over $N$ samples, without matrix multiplications).
    GA-SLM is the most expensive at inference time due to its online
    optimization.  The DNN-based approaches (PRNet, NN-SCF, ESLM-AE,
    DL-AE) have moderate inference cost dominated by the FC layer
    multiplications~\cite{nguyen2022deep}.
\end{keyinsight}


% ----------------------------------------------------------------------------
\section{Side Information and Blind Detection}
\label{sec:comparison_si}
% ----------------------------------------------------------------------------

\begin{table}[H]
    \centering
    \caption{Side information requirements.
             Data from~\cite{kim2017novel,zou2021novel,ali2016new,hao2019extended,nguyen2022deep,alnaseri2025deep}.}
    \label{tab:si_comparison}
    \begin{tabular}{@{} l c l @{}}
        \toprule
        \textbf{Approach} & \textbf{SI Bits} & \textbf{How SI is Avoided} \\
        \midrule
        Classical SLM & $\lceil\log_2 U\rceil$ & Not avoided \\
        PRNet & 0 & Joint encoder--decoder training \\
        NN-SCF & 0 & Tx/Rx modules share trained params \\
        GA-SLM & $\lceil\log_2 P\rceil$ & Not avoided \\
        ESLM-AE & 0 & AE decoder learns phase reversal \\
        PR-DUN & 0 & No phase rotation (additive $\bm{c}$) \\
        DL-AE & 0 & Joint encoder--decoder training \\
        \bottomrule
    \end{tabular}
\end{table}

Five of the six ML approaches eliminate side information entirely.
Only GA-SLM retains SI because it preserves the classical SLM framework
without a learned receiver~\cite{ali2016new}.


% ----------------------------------------------------------------------------
\section{BER Impact}
\label{sec:comparison_ber}
% ----------------------------------------------------------------------------

\begin{table}[H]
    \centering
    \caption{BER impact of each approach.
             Data from~\cite{kim2017novel,zou2021novel,ali2016new,hao2019extended,nguyen2022deep,alnaseri2025deep}.}
    \label{tab:ber_comparison}
    \begin{tabular}{@{} l c l @{}}
        \toprule
        \textbf{Approach} & \textbf{BER Degradation} & \textbf{Mechanism} \\
        \midrule
        Classical SLM & None (if SI correct) & Distortion-free \\
        PRNet & Minimal & Decoder compensates \\
        NN-SCF & Tunable via $\alpha$ & Trade-off parameter \\
        GA-SLM & None & Distortion-free (phase rotation) \\
        ESLM-AE & Minimal & Decoder compensates \\
        PR-DUN & Minimal & Power constraint limits distortion \\
        DL-AE & Minimal & Decoder compensates \\
        \bottomrule
    \end{tabular}
\end{table}


% ----------------------------------------------------------------------------
\section{Key Trade-off Patterns}
\label{sec:tradeoff_patterns}
% ----------------------------------------------------------------------------

Several clear patterns emerge from the comparison:

\paragraph{Trade-off 1: Complexity vs.\ Flexibility.}
Neural network approaches (PRNet, NN-SCF, ESLM-AE, DL-AE) have
\emph{fixed} inference cost regardless of the desired PAPR level, but
cannot adapt to different system configurations without retraining.
GA-SLM has \emph{variable} cost (adjustable via $P$ and $G$) and
naturally adapts to each symbol, but is computationally
expensive~\cite{ali2016new,kim2017novel}.

\paragraph{Trade-off 2: Interpretability vs.\ Performance.}
PR-DUN offers the highest interpretability (each layer = one PGD step)
with the fewest parameters, but introduces signal distortion.
PRNet and DL-AE offer strong performance with no SI but are
``black boxes''---it is hard to understand \emph{what} the network
has learned~\cite{nguyen2022deep,kim2017novel}.

\paragraph{Trade-off 3: Distortion-Free vs.\ SI-Free.}
The distortion-free approaches (GA-SLM, ESLM-AE) preserve the exact
constellation points, ensuring zero theoretical BER degradation.
The distortion-based approaches (PRNet, NN-SCF, PR-DUN, DL-AE) may
introduce small signal distortion but eliminate side information.
Only ESLM-AE achieves both distortion-free operation \emph{and}
SI elimination~\cite{hao2019extended}.

\begin{keyinsight}
    ESLM-AE is the only approach that simultaneously achieves:
    (1)~distortion-free operation (preserves SLM phase-rotation structure),
    (2)~no side information (autoencoder learns the decoder),
    (3)~fast inference (single forward pass + one IFFT).
    This makes it arguably the most balanced approach,
    though it has only been demonstrated for DCO-OFDM
    (not RF OFDM)~\cite{hao2019extended}.
\end{keyinsight}


% ----------------------------------------------------------------------------
\section{Practical Recommendations}
\label{sec:practical_recommendations}
% ----------------------------------------------------------------------------

Based on the comparative analysis above, the following recommendations
target engineers selecting a technique for a specific deployment scenario:

\paragraph{Scenario 1: Real-Time RF OFDM (e.g., 5G base station).}
\textbf{Recommended:} PR-DUN~\cite{nguyen2022deep}.
Its $\mathcal{O}(N)$ per-layer cost, minimal trainable parameters ($\sim 20$),
and built-in power constraint make it the most practical choice for
latency-sensitive RF systems.  The slight signal distortion is an acceptable
trade-off given the dramatic complexity reduction over classical ICF.

\paragraph{Scenario 2: Visible Light Communication (VLC / Li-Fi).}
\textbf{Recommended:} ESLM-AE~\cite{hao2019extended}.
It is specifically designed for DCO-OFDM, preserves the SLM structure
(distortion-free), eliminates side information, and has been validated
under realistic VLC channels (LOS, Rician, DOW).

\paragraph{Scenario 3: Long-Haul Coherent Optical Fiber.}
\textbf{Recommended:} DL-AE~\cite{alnaseri2025deep}.
It is the only approach tailored to the CO-OFDM domain, where Kerr
nonlinearity makes PAPR reduction particularly impactful.  The autoencoder
additionally improves the received constellation quality, a secondary benefit
unique to the nonlinear fiber channel.

\paragraph{Scenario 4: Research Prototyping and Exploration.}
\textbf{Recommended:} PRNet~\cite{kim2017novel} as a starting point.
Its simple autoencoder architecture, well-documented training procedure,
and reproducible parameters make it ideal for researchers entering the
field.  For ablation studies or algorithm design, GA-SLM~\cite{ali2016new}
provides interpretable per-symbol optimization that can serve as an
upper bound on achievable PAPR reduction.

\begin{keyinsight}
    No single approach dominates across all scenarios.  The ``best'' choice
    depends on: (a)~the OFDM domain (RF, optical intensity, coherent optical),
    (b)~the latency budget, (c)~whether distortion-free operation is required,
    and (d)~whether a custom receiver (decoder) can be deployed alongside the
    transmitter (encoder).
\end{keyinsight}
